\documentclass[11pt]{article}
\usepackage[a4paper,margin=25mm]{geometry}
\usepackage{amsmath,amssymb,amsthm,hyperref}
\newtheorem{proposition}{Proposition}
\title{Full fairness does not force independence of prefix coordinates}
\author{Exact obstruction to a proposed rank argument; independently reviewed}
\date{30 September 2026}
\begin{document}\maketitle
This note concerns a possible route to a factorial lower bound for the
total number of faces. It shows that full permutation fairness does not
force the raw ordered-prefix coordinates to be linearly independent.
It does not settle regularity of the endpoint map in the reduced Lie
algebra, whose tangent vectors involve conjugation by both prefixes and
their inverses.

For a finite word $T$, let $c_i$ be its total number of occurrences of
letter $i$. For a word $u=i_1\cdots i_d$ with distinct letters and a cut
$t$ in $T$, define
\[
 F_u(t)=\frac{\#\{\text{subsequences }i_1\cdots i_d
                    \text{ entirely before }t\}}
                  {c_{i_1}\cdots c_{i_d}}.
\]
Thus $F_u$ is a normalized prefix coordinate of the axis path of $T$.
One may either use cuts between letters or interpolate each individual
letter occurrence continuously; the conclusion below holds in both
conventions.

\begin{proposition}
For every $n\ge2$ there is a fully permutation-fair $n$-letter word $T$
and a distinguished letter $c$ such that, for each $0\le d\le n-1$,
the functions
\[
 \{F_{u c}: u\text{ is any ordered list of }d
                   \text{ distinct letters other than }c\}
\]
are all identical on every cut of $T$.
Consequently this family, which has $(n-1)!/(n-1-d)!$ members, has
linear span of dimension one whenever $d\ge1$.
\end{proposition}
\begin{proof}
Take a fair word $W$ on the $n-1$ old letters, with old counts $a_i$.
Choose a rational equal-weight quadrature $x_1,\ldots,x_K$ for the
uniform interval, exact through degree $n-1$, and a common denominator
$N$ of its nodes. The rational-design insertion proposition in
\texttt{main.tex} constructs a fair word $T$ by concatenating $N$
copies of $W$ and inserting one $c$ in gap $Nx_\ell$ for every node.
Repeated nodes give consecutive occurrences of $c$. All old letters
have total counts $Na_i$, and $c$ has total count $K$.

Before a $c$ in gap $j$, the restriction of the prefix to old letters
is $W^j$. Self-concatenation preserves fairness. Hence, for every
ordered list $u$ of $d$ distinct old letters, its normalized count in
that prefix is
\[
 \frac{j^d\prod_{i\in u}a_i/d!}{\prod_{i\in u}Na_i}
       =\frac{(j/N)^d}{d!}.
\]
The coordinate $F_{u c}$ increases only at occurrences of $c$. At the
occurrence corresponding to node $x_\ell$, its increment is therefore
$x_\ell^d/(K d!)$, independently of $u$. All these coordinates start
at zero, so at every cut
\[
 F_{u c}(t)=\frac1{K d!}
       \sum_{\substack{\ell:\text{the corresponding }c
                                  \text{ lies before }t}}x_\ell^d.
\]
This also proves the interpolated assertion. At the final cut the
common value is $1/(d+1)!$ by quadrature exactness, so it is nonzero.
\end{proof}

The construction works with unequal old counts as well. Thus the
collapse holds across different old subsets, not just permutations of
one fixed subset. In particular $F_{abc}=F_{bac}$ for distinct old
letters $a,b$, even in fully fair families with arbitrarily many dice.
A Gram matrix of these raw prefix functions is singular regardless of
how high a fairness strength is assumed. A valid lower-bound argument
using endpoint derivatives would need additional structure beyond
independence of these prefix coordinates; the proposition does not
assert that their equality produces a Lie-algebra tangent annihilator.
\end{document}
